% 【本文件写什么】impl 实现层：loongarch64
% - 定位：LoongArch 页表 / 架构相关 impl。
% - crate 路径：os/components/wateros-mm/mm-impl/impl-loongarch64/src/
% - 如何实现对应 api-v0 trait：关键类型、状态机、数据结构
% - 算法或硬件交互流程，可用伪代码/流程图
% - 本 impl 启用的 Cargo [features] 与 arch feature，impl-riscv64 等
% - 与其它 impl 变体的差异、切换 feature 时的影响
% - 性能/内存假设与已知限制
% 【事实来源】
% - 上述 crate 源码与 Cargo.toml
% - os/feature-tree.txt 中指向本 impl 的 feature 链

\subsubsection{impl-loongarch64}

LoongArch64 页表实现与 Sv39 在语义层对齐，包括三级 4\,KiB 叶子页、内核恒等映射、ELF\code{PT\_LOAD}、惰性 mmap、brk/munmap/mprotect、mremap 子集、fork COW、软件 walk 用户拷贝，硬件编码不同。

地址空间 token：\code{AddressSpaceOps::satp\_value}返回\code{root\_ppn * PAGE\_SIZE}，供 CSR.PGDL 直接写入；\code{LoadedElf::satp}字段名保留历史命名，实际内容为 PGDL 值。

内核 bring-up：\code{kernel\_global::init}在 S 态安装 PGDL 后映射 QEMU virt RAM 与 MMIO，含 VirtIO PCI BAR 窗口。PGDL 切换后用帧池内真实 RAM 帧做访存探针，避免低地址 VA 受 PGDL/PGDH 选择规则影响。

fork 与 COW：\code{fork\_user\_aspace}与\code{handle\_cow\_fault}接口与 Sv39 相同，内部\code{LoongArch64AddressSpace}实现\code{fork\_cow}与\code{handle\_cow\_page}，同样依赖\code{frame\_inc\_ref}/\code{frame\_dealloc\_result}管理共享帧。

trap 保留区：与 RISC-V 依赖链接脚本符号不同，LoongArch 路径结合 RAM 区间判断用户栈增长是否越界，具体见\code{user\_heap\_mmap}与缺页处理。

用户访问：\code{LoongArch64UserMemoryOps}实现\code{UserMemoryOps}，缺页写路径同样先尝试 COW 修复。无\code{debug\_probe\_user\_virt}诊断导出。

聚合层通过\code{cfg(feature = "impl-loongarch64")}选择\code{ActiveUserMemoryOps}别名与\code{kernel\_mm}转发；与\code{impl-sv39}互斥编译。

\begin{rustcode}
fn satp_value(&self) -> usize { self.root.0 * PAGE_SIZE }
\end{rustcode}
